\centerline{\bf \Huge ДЗ. Числовые таблицы.} 

\centerline{\bf \Huge Сложный перебор и чётность.} 

\centerline{\bf \Huge }

\begin{flushright}
        \scriptsize{\textbf{Откуда такие деньги?!}}
\end{flushright}

\problem{\textbf{1}} Какое минимальное количество клеток в квадрате 7$\times$7 нужно закрасить, чтобы в каждом прямоугольнике 2$\times$3 было по одной закрашенной клетке? Сможете доказать оценку?

\problem{\textbf{2}} Можно ли заполнить таблицу 5$\times$5 числами от 1 до 25 так, чтобы сумма чисел в каждом столбце была одинаковой?

\problem{\textbf{3}} Сколько шестизначных чисел может составить НеУченик из цифр 0, 2 и 5, если цифры могут повторяться?

\problem{\textbf{4}} Сколько чётных чисел от 13 до 2025 имеют не больше двух одинаковых цифр?

\problem{\textbf{5}} У АА есть огромная копилка, в которую он закинул какое-то количество пятитысячных купюр. Ученики закинули в эту же копилку стодолларовые купюры, которых было на 52 больше, чем пятитысячных купюр АА. Может ли в копилке оказаться ровно 2025 купюр?

\problem{\textbf{6}} АА загадал 2025 различных чисел. Ученик утверждает, что, если сумма любых пятидесяти двух чисел из них нечётна, то все 2025 чисел нечётны. НеУченик утверждает, что, если сумма любых пятидесяти трёх чисел из них нечётна, то все 2025 чисел нечётны. Кто прав, а кто не прав? Объясните почему?